% ============================================================================
% trajes.tex  —  DISFRACES/accesorios intercambiables del capibara (obra propia, CC0)
% ----------------------------------------------------------------------------
% Overlays PLANOS y RECOLOREABLES dibujados en LAS MISMAS COORDENADAS que
% capibara_kawaii.tex (y-arriba). Se ponen SOBRE el capibara sin taparle la cara
% (zona segura de la cara: ojos en (+-0.54,0.50) tope y~0.85; nariz (0,0); boca
% (0,-0.30); cachetes (+-1.00,-0.02)).
%
% Requiere: % ============================================================================
% _estilo_interes.tex  —  Estilo compartido del "mundo del capibara" (CC0)
% ----------------------------------------------------------------------------
% Paleta plana + estilo de linea + FRAME canonico para trajes/accesorios que se
% dibujan EN LAS MISMAS COORDENADAS que el capibara base
% (tikzlib/interes/capibara_kawaii.tex). Todo obra propia, CC0.
%
% Requiere: \usepackage{tikz}. NO usa texto/fuentes -> apto SVGMobject de Manim.
%
% FRAME CANONICO DE REGISTRO (\capiframe):
%   Todos los trajes/accesorios y el propio capibara se exportan con ESTE mismo
%   \useasboundingbox. Asi los SVG comparten viewBox y en Manim se pueden cargar
%   con  SVGMobject(..., should_center=False)  y quedan PERFECTAMENTE alineados
%   sobre el capibara sin tocar nada (registro automatico).
%   Con should_center=True (por defecto) el frame se ignora (Manim recorta al
%   contenido visible), asi que NO rompe usos existentes del capibara.
% ============================================================================

% --- linework del capibara (idempotente con capibara_kawaii.tex) ---
\providecolor{capicontorno}{HTML}{6E4E34}  % linea (marron calido suave)
\tikzset{
  capilinea/.style   ={draw=capicontorno, line width=2.1pt,
                       line cap=round, line join=round},
  capilineaf/.style  ={draw=capicontorno, line width=1.6pt,
                       line cap=round, line join=round}, % linea fina
}

% --- paleta de trajes/accesorios (plana, tierna) ---
\providecolor{vueltiaobase}{HTML}{F2E6C6}  % cana flecha (crema claro)
\providecolor{vueltiaosom} {HTML}{E4D2A6}  % sombra del ala
\providecolor{vueltiaopin} {HTML}{3B332B}  % "pintao" (bandas oscuras)

\providecolor{banqcorbata} {HTML}{9C3B4C}  % corbatin vino
\providecolor{banqcorbdark}{HTML}{7E2E3D}  % corbatin sombra
\providecolor{oro}         {HTML}{D2A63C}  % oro (monoculo/detalles)
\providecolor{orodark}     {HTML}{B4862A}

\providecolor{comerdelan}  {HTML}{5E86A6}  % delantal azul mercado
\providecolor{comerdeldark}{HTML}{4B6E8C}
\providecolor{comergorra}  {HTML}{CF5B45}  % gorra roja
\providecolor{comergordark}{HTML}{B24734}

\providecolor{acadboard}   {HTML}{343747}  % birrete azul-noche
\providecolor{acadbdark}   {HTML}{262838}
\providecolor{acadgafas}   {HTML}{3B332B}  % marco de gafas

\providecolor{ahorrosa}    {HTML}{F2A6BC}  % cerdito rosado
\providecolor{ahorrodark}  {HTML}{E086A0}
\providecolor{ahorrogorro} {HTML}{D9A441}  % gorro mostaza

\providecolor{naranjafill} {HTML}{F59E2E}  % naranja
\providecolor{naranjadark} {HTML}{E07C1A}  % naranja sombra
\providecolor{naranjacont} {HTML}{9A5A1E}  % contorno naranja
\providecolor{hojaverde}   {HTML}{6BA83C}  % hoja
\providecolor{hojadark}    {HTML}{4E8A2E}
\providecolor{canastamimb} {HTML}{C8925A}  % mimbre canasta
\providecolor{canastadark} {HTML}{A9743F}

% --- fondo/escena ---
\providecolor{escpared}    {HTML}{FCF4E2}  % pared crema (combina con el video)
\providecolor{escpiso}     {HTML}{EFE1C4}  % piso
\providecolor{esczocalo}   {HTML}{E4D2A6}
\providecolor{banqcol}     {HTML}{D9E3EC}  % columnas banco
\providecolor{banqcoldark} {HTML}{BFCEDC}
\providecolor{banqarco}    {HTML}{C7D6E4}
\providecolor{merctoldoA}  {HTML}{D9694E}  % toldo rojo
\providecolor{merctoldoB}  {HTML}{FBF1DC}  % toldo crema
\providecolor{mercmadera}  {HTML}{C8925A}  % madera del puesto
\providecolor{mercmaddark} {HTML}{A9743F}
\providecolor{fincacielo}  {HTML}{CDE7F0}  % cielo
\providecolor{fincapasto}  {HTML}{9FCB6B}  % pasto
\providecolor{fincapastod} {HTML}{7FB050}
\providecolor{fincasol}    {HTML}{F6C64B}  % sol
\providecolor{fincamata}   {HTML}{6BA83C}

% FRAME CANONICO (capibara + trajes) — simetrico en x; contiene sombreros altos
% (hasta y~2.05) y la alcancia lateral (hasta x~2.65).
\newcommand{\capiframe}{\useasboundingbox (-2.75,-2.10) rectangle (2.75,2.25);}

%
% Uso:  \trajevueltiao  \trajebanquero  \trajecomerciante  \trajeacademico
%       \trajeahorrador
% Recolorear: \colorlet{vueltiaobase}{...} etc. (ver _estilo_interes.tex).
%
% NOTAS DE POSICION/ESCALA (relativas al capibara base):
%   - Todos se dibujan ya colocados; en Manim carga capibara y traje con
%     SVGMobject(..., should_center=False) y quedan alineados (mismo \capiframe).
%   - Alternativa manual: cada traje trae su ancla en el comentario de su macro.
% ============================================================================

% ---------------------------------------------------------------------------
% 1) SOMBRERO VUELTIAO  (para "explicar en general" / toque colombiano)
%    Ala ancha + copa baja; cana flecha crema con bandas "pintao" oscuras.
%    Descansa sobre la cabeza, por ENCIMA de los ojos. Centro-x = capibara.
%    Ancla: centro del ala en (0,1.30); copa sube a y~2.02.
% ---------------------------------------------------------------------------
\newcommand{\trajevueltiao}{%
  \begin{scope}
    % ala (elipse ancha)
    \path[capilinea, fill=vueltiaobase] (0,1.30) ellipse[x radius=1.78, y radius=0.34];
    \path[fill=vueltiaosom] (0,1.16) ellipse[x radius=1.66, y radius=0.20];
    % copa (cupula redondeada)
    \path[capilinea, fill=vueltiaobase, rounded corners=0.42cm]
        (-0.84,1.28) rectangle (0.84,2.00);
    % bandas "pintao" (dos franjas oscuras en la copa)
    \path[fill=vueltiaopin, rounded corners=0.10cm] (-0.80,1.44) rectangle (0.80,1.60);
    \path[fill=vueltiaopin, rounded corners=0.10cm] (-0.72,1.74) rectangle (0.72,1.86);
    % franja del ala (acento)
    \path[fill=vueltiaopin] (0,1.30) ellipse[x radius=1.70, y radius=0.075];
    \path[fill=vueltiaobase] (0,1.335) ellipse[x radius=1.70, y radius=0.05];
  \end{scope}
}

% ---------------------------------------------------------------------------
% 2) BANQUERO  (para "el banco"): corbatin vino + monoculo dorado con cadenita.
%    Corbatin a la altura del cuello (0,-0.90). Monoculo rodea el ojo DERECHO.
% ---------------------------------------------------------------------------
\newcommand{\trajebanquero}{%
  \begin{scope}
    % --- corbatin (bow tie) ---
    \path[capilineaf, fill=banqcorbata]
        (-0.05,-0.90) -- (-0.52,-0.66) -- (-0.52,-1.14) -- cycle;
    \path[capilineaf, fill=banqcorbata]
        ( 0.05,-0.90) -- ( 0.52,-0.66) -- ( 0.52,-1.14) -- cycle;
    \path[fill=banqcorbdark] (-0.44,-0.90) ellipse[x radius=0.05, y radius=0.16];
    \path[fill=banqcorbdark] ( 0.44,-0.90) ellipse[x radius=0.05, y radius=0.16];
    \path[capilineaf, fill=banqcorbdark, rounded corners=0.05cm]
        (-0.11,-1.02) rectangle (0.11,-0.78);
    % --- monoculo sobre el ojo derecho ---
    \draw[capilinea, line width=2.4pt, draw=oro] (0.54,0.50) circle[radius=0.48];
    \draw[capilineaf, draw=oro] (0.90,0.14) .. controls (1.02,-0.20) .. (0.86,-0.52);
    \path[fill=orodark] (0.86,-0.52) circle[radius=0.06];
  \end{scope}
}

% ---------------------------------------------------------------------------
% 3) COMERCIANTE  (para "el mercado"): delantal azul + gorra roja con visera.
%    Delantal sobre la barriga (no tapa la cara). Gorra sobre la cabeza.
% ---------------------------------------------------------------------------
\newcommand{\trajecomerciante}{%
  \begin{scope}
    % --- delantal ---
    % tirantes al cuello
    \draw[capilineaf, draw=comerdeldark] (-0.30,-0.60) -- (-0.14,-0.30);
    \draw[capilineaf, draw=comerdeldark] ( 0.30,-0.60) -- ( 0.14,-0.30);
    % peto (bib)
    \path[capilineaf, fill=comerdelan, rounded corners=0.12cm]
        (-0.56,-1.16) rectangle (0.56,-0.60);
    % falda del delantal
    \path[capilineaf, fill=comerdelan, rounded corners=0.20cm]
        (-1.06,-1.98) rectangle (1.06,-1.10);
    % bolsillo
    \path[capilineaf, fill=comerdeldark, rounded corners=0.06cm]
        (-0.42,-1.70) rectangle (0.42,-1.34);
    \draw[capilineaf, draw=comerdelan] (0,-1.34) -- (0,-1.70);
    % --- gorra roja con visera ---
    \path[capilinea, fill=comergorra] (0,1.46) ellipse[x radius=0.98, y radius=0.60];
    \path[fill=comergordark] (0,1.10) ellipse[x radius=0.94, y radius=0.22];
    % visera (mira al frente-abajo)
    \path[capilinea, fill=comergorra] (0,1.02) ellipse[x radius=0.70, y radius=0.20];
    % botoncito
    \path[fill=comergordark] (0,2.02) circle[radius=0.08];
  \end{scope}
}

% ---------------------------------------------------------------------------
% 4) ACADEMICO  (para "explicar"): birrete (mortarboard) + gafas redondas.
%    Birrete sobre la cabeza; gafas rodean AMBOS ojos.
% ---------------------------------------------------------------------------
\newcommand{\trajeacademico}{%
  \begin{scope}
    % --- birrete ---
    % banda que apoya en la cabeza
    \path[capilinea, fill=acadboard, rounded corners=0.10cm]
        (-0.66,1.30) rectangle (0.66,1.66);
    % tabla (rombo aplanado, vista frontal)
    \path[capilinea, fill=acadboard]
        (0,2.04) -- (1.22,1.66) -- (0,1.28) -- (-1.22,1.66) -- cycle;
    \path[fill=acadbdark] (0,1.66) circle[radius=0.07]; % boton central
    % borla (tassel)
    \draw[capilineaf, draw=oro] (1.22,1.66) .. controls (1.28,1.30) .. (1.16,0.98);
    \path[fill=oro] (1.16,0.98) ellipse[x radius=0.10, y radius=0.18];
    \path[fill=orodark] (1.16,1.06) ellipse[x radius=0.10, y radius=0.05];
    % --- gafas redondas ---
    \draw[capilinea, line width=2.2pt, draw=acadgafas] (-0.54,0.50) circle[radius=0.46];
    \draw[capilinea, line width=2.2pt, draw=acadgafas] ( 0.54,0.50) circle[radius=0.46];
    \draw[capilinea, line width=2.2pt, draw=acadgafas] (-0.08,0.52) -- (0.08,0.52);
    \draw[capilineaf, draw=acadgafas] (-1.00,0.54) -- (-1.26,0.66);
    \draw[capilineaf, draw=acadgafas] ( 1.00,0.54) -- ( 1.26,0.66);
  \end{scope}
}

% ---------------------------------------------------------------------------
% 5) AHORRADOR  (para "ahorrar"): gorrito mostaza + alcancia-cerdito AL LADO
%    (con una naranja-moneda entrando por la ranura). El cerdito va a la derecha
%    sobre el piso; el gorrito es opcional (usalo o el cerdito, o ambos).
% ---------------------------------------------------------------------------
\newcommand{\trajeahorrador}{%
  \begin{scope}
    % --- gorrito tejido (dome + vuelta + pompon) ---
    \path[capilinea, fill=ahorrogorro] (0,1.50) ellipse[x radius=0.96, y radius=0.62];
    \path[capilinea, fill=ahorrogorro, rounded corners=0.12cm]
        (-0.98,0.98) rectangle (0.98,1.22);
    \path[capilinea, fill=ahorrogorro] (0,2.14) circle[radius=0.18]; % pompon
    % rayitas del tejido
    \draw[capilineaf, draw=orodark] (-0.30,1.10) -- (-0.30,1.90);
    \draw[capilineaf, draw=orodark] ( 0.30,1.10) -- ( 0.30,1.90);
    % --- alcancia cerdito (al lado derecho, sobre el piso) ---
    \begin{scope}
      % patitas
      \path[fill=ahorrodark] (1.62,-1.72) rectangle (1.80,-1.52);
      \path[fill=ahorrodark] (2.14,-1.72) rectangle (2.32,-1.52);
      % cuerpo
      \path[capilinea, fill=ahorrosa] (1.98,-1.14) ellipse[x radius=0.64, y radius=0.50];
      % hocico
      \path[capilinea, fill=ahorrodark] (2.56,-1.16) ellipse[x radius=0.20, y radius=0.24];
      \path[fill=vueltiaopin] (2.56,-1.10) circle[radius=0.045];
      \path[fill=vueltiaopin] (2.56,-1.24) circle[radius=0.045];
      % orejita
      \path[capilinea, fill=ahorrosa] (2.30,-0.66) -- (2.52,-0.78) -- (2.30,-0.90) -- cycle;
      % ojo
      \path[fill=capicontorno] (2.28,-1.02) circle[radius=0.055];
      % ranura + naranja-moneda entrando
      \path[fill=vueltiaopin, rounded corners=0.03cm] (1.80,-0.70) rectangle (2.16,-0.62);
      \path[capilineaf, fill=naranjafill] (1.98,-0.34) circle[radius=0.20];
      \path[fill=hojaverde] (2.06,-0.16) ellipse[x radius=0.08, y radius=0.04];
      % colita
      \draw[capilineaf, draw=ahorrodark] (1.34,-1.10) .. controls (1.18,-1.02) and (1.22,-1.28) .. (1.10,-1.20);
    \end{scope}
  \end{scope}
}
